\documentclass[10pt,a4paper]{amsart}
\usepackage{amsmath,amssymb,amsthm,graphicx}
\usepackage{verbatim}
\newenvironment{reference}{\comment}{\endcomment}
%\newenvironment{reference}{}{}
\newenvironment{slogan}{\comment}{\endcomment}
\newenvironment{history}{\comment}{\endcomment}
% For commutative diagrams we use Xy-pic
\usepackage[all]{xy}

% We use 2cell for 2-commutative diagrams.
\xyoption{2cell}
% Theorem environments.
%
\theoremstyle{plain}
\newtheorem{theorem}[subsection]{Theorem}
\newtheorem{proposition}[subsection]{Proposition}
\newtheorem{lemma}[subsection]{Lemma}

\theoremstyle{definition}
\newtheorem{definition}[subsection]{Definition}
\newtheorem{example}[subsection]{Example}
\newtheorem{exercise}[subsection]{Exercise}
\newtheorem{situation}[subsection]{Situation}

\theoremstyle{remark}
\newtheorem{remark}[subsection]{Remark}
\newtheorem{remarks}[subsection]{Remarks}

\numberwithin{equation}{subsection}

% Macros
%
\def\lim{\mathop{\mathrm{lim}}\nolimits}
\def\colim{\mathop{\mathrm{colim}}\nolimits}
\def\Spec{\mathop{\mathrm{Spec}}}
\def\Hom{\mathop{\mathrm{Hom}}\nolimits}
\def\Ext{\mathop{\mathrm{Ext}}\nolimits}
\def\SheafHom{\mathop{\mathcal{H}\!\mathit{om}}\nolimits}
\def\SheafExt{\mathop{\mathcal{E}\!\mathit{xt}}\nolimits}
\def\Sch{\mathit{Sch}}
\def\Mor{\mathop{\mathrm{Mor}}\nolimits}
\def\Ob{\mathop{\mathrm{Ob}}\nolimits}
\def\Sh{\mathop{\mathit{Sh}}\nolimits}
\def\NL{\mathop{N\!L}\nolimits}
\def\CH{\mathop{\mathrm{CH}}\nolimits}
\def\proetale{{pro\text{-}\acute{e}tale}}
\def\etale{{\acute{e}tale}}
\def\QCoh{\mathit{QCoh}}
\def\Ker{\mathop{\mathrm{Ker}}}
\def\Im{\mathop{\mathrm{Im}}}
\def\Coker{\mathop{\mathrm{Coker}}}
\def\Coim{\mathop{\mathrm{Coim}}}

% Boxtimes
%
\DeclareMathSymbol{\boxtimes}{\mathbin}{AMSa}{"02}

%
% Macros for moduli stacks/spaces
%
\def\QCohstack{\mathcal{QC}\!\mathit{oh}}
\def\Cohstack{\mathcal{C}\!\mathit{oh}}
\def\Spacesstack{\mathcal{S}\!\mathit{paces}}
\def\Quotfunctor{\mathrm{Quot}}
\def\Hilbfunctor{\mathrm{Hilb}}
\def\Curvesstack{\mathcal{C}\!\mathit{urves}}
\def\Polarizedstack{\mathcal{P}\!\mathit{olarized}}
\def\Complexesstack{\mathcal{C}\!\mathit{omplexes}}
% \Pic is the operator that assigns to X its picard group, usage \Pic(X)
% \Picardstack_{X/B} denotes the Picard stack of X over B
% \Picardfunctor_{X/B} denotes the Picard functor of X over B
\def\Pic{\mathop{\mathrm{Pic}}\nolimits}
\def\Picardstack{\mathcal{P}\!\mathit{ic}}
\def\Picardfunctor{\mathrm{Pic}}
\def\Deformationcategory{\mathcal{D}\!\mathit{ef}}
\usepackage{cleveref}
\crefname{originallemma}{original lemma}{original lemmas}
\Crefname{originallemma}{Original lemma}{Original lemmas}
\makeatletter
\newlabel{algebra-lemma-completion-tensor}{{00MA}{source}{}{sourcepointer-1}{}}
\newlabel{algebra-lemma-completion-tensor@cref}{{[originallemma][0][]00MA}{[1][0][]source}{}{}{}}
\makeatother
\crefname{originallemma}{original lemma}{original lemmas}
\Crefname{originallemma}{Original lemma}{Original lemmas}
\makeatletter
\newlabel{lemma-proper-over-affine-cohomology-finite}{{08AS}{source}{}{sourcepointer-2}{}}
\newlabel{lemma-proper-over-affine-cohomology-finite@cref}{{[originallemma][0][]08AS}{[1][0][]source}{}{}{}}
\makeatother
\crefname{originaltheorem}{original theorem}{original theorems}
\Crefname{originaltheorem}{Original theorem}{Original theorems}
\makeatletter
\newlabel{theorem-formal-functions}{{08AZ}{source}{}{sourcepointer-3}{}}
\newlabel{theorem-formal-functions@cref}{{[originaltheorem][0][]08AZ}{[1][0][]source}{}{}{}}
\makeatother
\title{The theorem on formal functions over a complete Noetherian base, with the completion on the left hand side redundant}
\author{The Stacks Project Authors}
\date{stacks-project:spaces-cohomology.tex:local-snapshot-20261012; GNU FDL 1.2 or later}
\begin{document}
\maketitle
\noindent\textit{Editorial note.} All delivered mathematical text is the author's own, copied from the pinned Stacks Project chapter file spaces-cohomology.tex; no statement or proof step was rewritten, reordered or removed.\par\medskip
\noindent\textit{Editorial note.} Preamble provenance: the pinned chapter file opens with \textbackslash{}input\{preamble\} and carries no preamble text of its own, so the project's single Stacks-wide preamble staged at sources/stacks-spaces-cohomology/preamble.tex is excerpted verbatim as three spans, lines 8-12 (verbatim package and the project's own reference/slogan/history display environments), lines 14-18 (Xy-pic setup) and lines 154-217 (theorem declarations and mathematical macros). The wrapper class and the amsmath/amssymb/amsthm/graphicx base packages are supplied by the assembler. No macro was invented or redefined.\par\medskip
\noindent\textit{Editorial note.} Same-chapter pointers are located at their \textbackslash{}label lines inside the pinned chapter file; cross-chapter pointers are located in the project's own tag list staged in the snapshot (sources/stacks-chow/tags.txt), at the verbatim lines carrying their canonical tags.\par\medskip
\noindent\textit{Allowed standard prerequisites (editorial).}\begin{itemize}
\item The theorem on formal functions. Same chapter spaces-cohomology.tex, Theorem \textbackslash{}ref\{theorem-formal-functions\} (cited at proof line 4494). It is the general statement whose completion of the cohomology modules is being removed in the case of a complete base ring.
\item Finiteness of the cohomology of a coherent sheaf on a proper space over a Noetherian affine base. Same chapter spaces-cohomology.tex, label lemma-proper-over-affine-cohomology-finite (invoked at proof line 4497). It makes \$H\textasciicircum{}p(X, \textbackslash{}mathcal\{F\})\$ a finite \$A\$-module, hence complete for the \$I\$-adic topology.
\item A finite module over a Noetherian complete ring is complete. Algebra, label algebra-lemma-completion-tensor (invoked at proof line 4498). It shows that the completion of the left hand side is not necessary, which yields the displayed identification.
\end{itemize}\medskip
\ifdefined\hypertarget\hypertarget{sourcepointer-1}{}\fi
\noindent\textit{Source reference.} Original lemma 00MA: sources/stacks-chow/tags.txt, source lines 816--816. Navigation pointer for \textbackslash{}ref\{algebra-lemma-completion-tensor\} used inside the counted span (line 4498); the argument that consumes it is retained in the span.\par\medskip
\ifdefined\hypertarget\hypertarget{sourcepointer-2}{}\fi
\noindent\textit{Source reference.} Original lemma 08AS: sources/stacks-spaces-cohomology/pinned-original-source.tex, source lines 4012--4012. Navigation pointer for \textbackslash{}ref\{lemma-proper-over-affine-cohomology-finite\} used inside the counted span (line 4497); the argument that consumes it is retained in the span.\par\medskip
\ifdefined\hypertarget\hypertarget{sourcepointer-3}{}\fi
\noindent\textit{Source reference.} Original theorem 08AZ: sources/stacks-spaces-cohomology/pinned-original-source.tex, source lines 4425--4425. Navigation pointer for \textbackslash{}ref\{theorem-formal-functions\} used inside the counted span (line 4494); the argument that consumes it is retained in the span.\par\medskip
\begin{lemma}
\label{lemma-spell-out-theorem-formal-functions}
Let $A$ be a ring. Let $I \subset A$ be an ideal. Assume $A$ is
Noetherian and complete with respect to $I$.
Let $f : X \to \Spec(A)$ be a proper morphism of algebraic spaces.
Let $\mathcal{F}$ be a coherent sheaf on $X$. Then
$$
H^p(X, \mathcal{F}) = \lim_n H^p(X, \mathcal{F}/I^n\mathcal{F})
$$
for all $p \geq 0$.
\end{lemma}

\begin{proof}
This is a reformulation of the theorem on formal functions
(Theorem \ref{theorem-formal-functions}) in the
case of a complete Noetherian base ring. Namely, in this case the
$A$-module $H^p(X, \mathcal{F})$ is finite
(Lemma \ref{lemma-proper-over-affine-cohomology-finite}) hence
$I$-adically complete (Algebra, Lemma \ref{algebra-lemma-completion-tensor})
and we see that completion on the left hand side is not necessary.
\end{proof}

\end{document}
